\documentclass[12pt]{article}
\usepackage[utf8]{inputenc}
\usepackage[T1]{fontenc}
\usepackage{hyperref}
\title{Walser Language Resources for K-12 Curriculum Development}
\author{Compiled by Claude Sonnet 5 (claude-sonnet-5) for Andr\'as Kornai}
\date{}
\begin{document}
\maketitle

\section{Introduction}
Walser (ISO 639-3 \texttt{wae}) is a group of Alemannic German varieties spoken in Switzerland, Italy, Austria and Liechtenstein. It is endangered. In Italy the Gressoney variety is called Titsch. The search turned up little material designed for schools, so this report is short.

\section{Primary Educational Resources}
No textbook or primer was found in the search results. One tourism page for the Aosta Valley says Titsch is taught in schools to young children and that the Walser Cultural Centre in Gressoney runs open Titsch courses. This is a secondary, non-specialist source, and the claim has not been verified against an official school source.
\begin{itemize}
\item Titsch: the German dialect from Gressoney (tourism page): \url{https://www.lovevda.it/en/database/10/tradition/gressoney-la-trinite/titsch-the-german-dialect-from-gressoney/447}
\end{itemize}

\section{Linguistic and Reference Materials}
\begin{itemize}
\item Wikipedia article on Walser German / Titsch: \url{https://Www.wikipedia.org/wiki/Titsch}
\item Overview page ``La langue walser'' (CEFAN, Universit\'e Laval): \url{https://axl.cefan.ulaval.ca/europe/walser.htm}
\item BnF catalogue category for Walser dictionaries and glossaries. No specific title was seen: \url{https://data.bnf.fr/12167970/walser__dialecte_/en.pdf}
\end{itemize}

\section{Religious and Community Resources}
The Walser Cultural Centre in Gressoney is mentioned in the tourism page above. No other community organisations or religious materials were found in the search results.

\section{Digital and Audio Resources}
No audio, video or software resources were found.

\section{Curriculum Development Context}
Walser is dispersed across several countries and has several local varieties, so materials would have to be built variety by variety. The only school-related lead is the Gressoney Titsch teaching mentioned above. It needs checking with the Walser Cultural Centre. Little else was found, and further searching in German and Italian is advisable.
\end{document}